
\documentclass[../../../physics-standard_questions_all.tex]{subfiles}

\begin{document}

真空中に$xy$平面を定め，$y$軸上の点$(0, \hspace{1.2mm} L)$に
電荷$Q\,(Q>0)$の点電荷Aを固定し，さらに$y$軸上の点$(0, \hspace{1.2mm} -L)$に
同じく電荷$Q$の点電荷Bを固定した．クーロンの法則の比例定数を$k_0$とし，
クーロン力以外の力の影響は考えない．
点$(x, \hspace{1.2mm} y)$における電場の$x$成分を$E_x(x, \hspace{1.2mm} y)$，
$y$成分を$E_y(x, \hspace{1.2mm} y)$，
電位（無限遠基準）を$\phi(x, \hspace{1.2mm} y)$と表す．

\begin{enumerate}

    \item $E_x(x, \hspace{1.2mm} 0)$，$E_y(x, \hspace{1.2mm} 0)$，
    $\phi(x, \hspace{1.2mm} 0)$を求めよ．

    \item 正の点電荷P（質量$m$，電荷$q$）を，無限遠から
    点$(\sqrt{3}L, \hspace{1.2mm} 0)$までゆっくりと運ぶ．
    この際に必要な仕事$W$を求めよ．

    \item Pが$x$軸上をなめらかに動けるものとする．
    Pに点$(\sqrt{3}L, \hspace{1.2mm} 0)$で，$x$軸負方向に初速$v_0$を与える場合，
    Pが$x$軸負方向の無限遠へ飛び去るための$v_0$の条件を求めよ．

    \item 負の点電荷N（質量$m$，電荷$-q$）を，無限遠から
    点$(\sqrt{3}L, \hspace{1.2mm} 0)$までゆっくりと運ぶ．
    この際に必要な仕事$W^*$を求めよ．

    \item Nが$x$軸上をなめらかに動けるものとする．
    Nに点$(\sqrt{3}L, \hspace{1.2mm} 0)$で，$x$軸正方向に初速$v_0$を与える場合，
    Nが$x$軸正方向の無限遠へ飛び去るための$v_0$の条件を求めよ．

    \item Nを原点付近の$x$軸上の点で静かに放したところ，Nの運動は単振動と見なせるものとなった．
    その周期を求めよ．ただし，微小量の2次の項は無視せよ．

\end{enumerate}

\end{document}


